\chapter{Waves on a surface}

\section{Objects that ride a moving surface}

Many tracked objects do not move freely through a volume. They travel on a surface: tracers on a
fluid interface, weather systems on a planet, features on a deforming membrane, layers of cells on a
developing organism. When that surface is close to a sphere it is modelled as an evolving sphere-like
surface, a closed surface $\mathcal{M}_t$ reached from the unit sphere by a radial map,
\[
x \;\mapsto\; \tilde\rho(t, x)\, x, \qquad x \in \mathcal{S}^2,
\]
so motion estimation on the surface is posed on the round sphere once the radius $\tilde\rho$ is
known. Lang and Scherzer set this up for optical flow and give the full construction; this chapter
follows theirs.

\section{Fitting the surface}

The radius is fitted to sample points on the surface, the approximate centres of the tracked
objects, after centring them on a fitted sphere. The fit minimises a data term plus a Sobolev
seminorm penalty,
\[
\mathcal{F}_\beta(\tilde\rho) \;=\; \|\tilde\rho - \tilde\rho^{\delta}\|^2_{L^2(\mathcal{S}^2)}
+ \beta\, |\tilde\rho|^2_{H^s(\mathcal{S}^2)},
\]
with the radius expanded in scalar spherical harmonics. The seminorm is diagonal in the harmonic
basis, weighting degree $n$ by $(n(n+1))^s$, so smoothness is a statement about how fast coefficients
decay with degree. Image values on the surface are read as the maximum intensity along a narrow
radial band about it.

\section{The flow in vector spherical harmonics}

Every tangent vector field on the sphere is expanded uniquely in tangential vector spherical
harmonics of two types,
\[
\tilde{\mathbf{y}}^{(2)}_{nj} = \lambda_n^{-1/2}\, \nabla_{\mathcal{S}^2} \tilde Y_{nj}, \qquad
\tilde{\mathbf{y}}^{(3)}_{nj} = \lambda_n^{-1/2}\, \nabla_{\mathcal{S}^2} \tilde Y_{nj} \times \tilde{\mathbf{N}},
\]
with $\lambda_n = n(n+1)$. The first type is a gradient field and carries the part of a flow that
converges or diverges; the second is a rotated gradient and carries the part that circulates. The
expansion is therefore a Helmholtz-Hodge decomposition already, which is why vector spherical
harmonics are called the natural basis for it: a flow on a sphere splits uniquely into an
irrotational part, a solenoidal part, and a harmonic part. There is no tangential harmonic of degree
zero, since by the hairy ball theorem no non-vanishing tangent field is constant over the sphere.

The optical flow on the surface minimises a data term, the squared residual of the generalised
optical flow equation, plus $\alpha$ times the squared $H^1$ norm of the field. Expanded in vector
harmonics up to a chosen degree, it becomes a linear system $(A + \alpha D)v = b$ solved iteratively
to a tolerance. The integrals are evaluated on a refined icosahedron with piecewise quadratic
interpolation and a quadrature rule per triangle, and in the published implementation those
integrals, not the linear solve, dominate the run time.

\section{What the harmonics buy a tracker}

The decomposition is ordered by scale. Degree 1 of the circulating type is a rigid turn of the whole
surface, the motion a turning view or a turning subject produces everywhere at once. Low degrees
carry the large-scale flows of the surface, the waves that move whole neighbourhoods of objects
together. What no low degree explains is the motion particular to one object. A tracker that removes
the low degrees before it links is linking on that remainder, which is the order the first chapter
asks for.
